\documentclass[11pt]{amsart}

\usepackage{amsmath,amssymb,amsthm,mathtools}
\usepackage{enumitem}
\usepackage[hidelinks]{hyperref}
\usepackage{microtype}

% PeAIce claim-discipline macros (input from manuscript preambles)
% Status tags for theorem environments and inline use

\newcommand{\PeAIceFormal}{\textsc{Formal}}
\newcommand{\PeAIceProposed}{\textsc{Proposed}}
\newcommand{\PeAIceOpen}{\textsc{Open}}
\newcommand{\PeAIceKnown}{\textsc{Known}}
\newcommand{\PeAIceNumerics}{\textsc{Numerics}}
\newcommand{\PeAIceClosedNeg}{\textsc{Closed-Negative}}
\newcommand{\PeAIceClosedPos}{\textsc{Closed-Positive}}
\newcommand{\PeAIceAnalogy}{\textsc{Structural Analogy}}

\newcommand{\statusbox}[1]{%
  \par\noindent\fbox{\parbox{0.95\linewidth}{\small\textbf{Status.} #1}}\par\medskip
}

\newcommand{\firewallbox}[1]{%
  \par\noindent\fbox{\parbox{0.95\linewidth}{\small\textbf{Firewall.} #1}}\par\medskip
}

\newcommand{\PeAIceOpenTargets}{%
  The {R}iemann hypothesis and the {C}oleman conjecture remain \PeAIceOpen.
  No claim in this note constitutes a proof of either statement.
}

\newcommand{\ZetaFirewall}{%
  The value $\zeta(0)=-\tfrac12$ is \emph{not} asserted as a term in the
  {K}akeya proof literature. Any regularization at $s=0$ is second-stage and
  conditional on a residual operator or counting object constructed after
  residual control.
}

\newcommand{\KNSFirewall}{%
  Typed incidence objects do not locate zeros of $\zeta$.
  Overlap measures are not placement measures on the critical line.
}


\theoremstyle{plain}
\newtheorem{theorem}{Theorem}[section]
\newtheorem{lemma}[theorem]{Lemma}
\newtheorem{proposition}[theorem]{Proposition}
\newtheorem{corollary}[theorem]{Corollary}

\theoremstyle{definition}
\newtheorem{definition}[theorem]{Definition}
\newtheorem{remark}[theorem]{Remark}
\newtheorem{conjecture}[theorem]{Conjecture}

\newcommand{\KB}{\mathrm{KB}}
\newcommand{\RH}{\mathrm{RH}}
\newcommand{\LH}{\mathrm{LH}}
\newcommand{\CC}{\mathrm{CC}}
\newcommand{\NsigT}{N(\sigma,T)}

\title[The Coleman Conjecture]{The Coleman Conjecture:\\
Kakeya geometry as a necessary precursor of Riemann critical-line control}
\author{Manuel Coleman}
\address{Love Labs LCA / PeAIce Research Program}
\email{coleman.manuel23@gmail.com}
\date{\today}

\begin{document}

\begin{abstract}
We introduce and rigorize the \emph{Coleman Conjecture}: the proposal that
Kakeya--Besicovitch geometric structure is a \emph{necessary precursor}---prior
in the dependency order---of any control of the nontrivial zeros of the Riemann
zeta function that reaches the critical line, rather than a sufficient condition
for the Riemann hypothesis.
We prove a packaging theorem (the \emph{sufficiency trap}): with the Kakeya set
conjecture in $\mathbb{R}^3$ now a theorem, the reading ``Kakeya$(\mathbb{R}^3)
\Rightarrow\RH$'' either proves $\RH$ or merely relabels the missing proof.
We record the established incidence $\to$ restriction/decoupling $\to$
large-values $\to$ zero-density tower, whose outputs are bounds (culminating in
Lindel\"of-type ceilings), not exact zero locations.
We then state two non-sufficient forms of the conjecture: an
\emph{invariant form} (existence of a faithful shared exponent/invariant
$\kappa$ whose extremal value is equivalent to $\RH$) and a
\emph{construction form} (any self-adjoint realization of the completed
xi-function by a $\zeta$-regularized determinant must be built on
Kakeya-type incidence data, and after the closed square-difference determinant
corridor must be prime-carrying).
Open problems, falsifiers, and claim-status tags are listed explicitly.
\PeAIceOpenTargets
\end{abstract}

\subjclass[2020]{11M26, 42B25, 42B20, 11N05}
\keywords{Coleman Conjecture, Riemann hypothesis, Kakeya set conjecture,
restriction and decoupling, zero-density estimates, Lindel\"of hypothesis,
Hilbert--P\'olya, prime-carrying operators}

\maketitle

\section{Introduction}

\statusbox{Program note ARX-005 under PeAIce claim discipline.
Sufficiency packaging is \PeAIceClosedNeg\ as a research packaging of $\RH$.
Necessary-precursor reading is \PeAIceProposed/\PeAIceOpen.
No proof of $\RH$ or of the Coleman Conjecture is claimed.}

The Kakeya--Besicovitch phenomenon exhibits directional completeness under
severe volume constraints: a compact set in $\mathbb{R}^n$ may contain a unit
segment in every direction while having arbitrarily small measure, and in
$\mathbb{R}^3$ every such set has Minkowski and Hausdorff dimension three
\cite{WangZahl2025Kakeya,GuthWangZahl2026Streamlined}.
The Riemann hypothesis ($\RH$) asserts that every nontrivial zero of
$\zeta(s)$ lies on the critical line $\mathrm{Re}(s)=\tfrac12$.

The \emph{Coleman Conjecture} ($\CC$), named in the PeAIce research program,
proposes a structural relation between these two worlds. In informal shorthand
one sometimes writes
\begin{equation}
  \label{eq:slogan}
  \KB \;\longrightarrow\; \mathrm{Re}(s)=\tfrac12.
\end{equation}
The content of \eqref{eq:slogan} depends entirely on the reading of the arrow.
This paper fixes a non-degenerate reading and discards a degenerate one.

\begin{definition}[Two readings of antecedence]
\label{def:readings}
\begin{enumerate}[label=(\alph*)]
  \item \textbf{Reading (S) --- sufficient condition.}
  Kakeya-type structure is the hypothesis of an implication that yields $\RH$.
  \item \textbf{Reading (N) --- necessary precursor.}
  Kakeya-type control is prior in the dependency order of any route that
  meaningfully constrains the zeros: a structural layer that successful methods
  pass through, not a package that by itself proves $\RH$.
\end{enumerate}
\end{definition}

\begin{conjecture}[Coleman Conjecture, necessary-precursor form]
\label{conj:cc-n}
\statusbox{\PeAIceProposed/\PeAIceOpen.}
Kakeya--Besicovitch incidence geometry (tube families, non-clustering, multi-scale
directional organization) is a necessary structural layer for any zero-facing
route that reaches exact critical-line control of $\zeta$, in the sense made
precise by the invariant form (Conjecture~\ref{conj:invariant}) and/or the
construction form (Conjecture~\ref{conj:construction}) below.
In particular, the conjecture does \emph{not} assert
\begin{equation}
  \label{eq:no-S}
  \text{Kakeya}(\mathbb{R}^3)\;\Longrightarrow\;\RH.
\end{equation}
\end{conjecture}

The paper is organized as follows.
Section~\ref{sec:trap} proves the sufficiency trap.
Section~\ref{sec:tower} records the dependency tower and the Lindel\"of ceiling.
Section~\ref{sec:invariant} states the invariant form and the faithfulness
obstruction.
Section~\ref{sec:construction} states the construction form and the closed
square-difference corridor.
Section~\ref{sec:open} lists open problems and falsifiers.

\begin{remark}[Claim tags]
Throughout we mark statements as \PeAIceKnown, \PeAIceFormal, \PeAIceProposed,
\PeAIceOpen, \PeAIceClosedNeg, or \PeAIceAnalogy, in the sense of an explicit
research-state grammar: definitions and deductions are not sold as theorems of
$\RH$, and open reverse implications are not collapsed into slogans.
\end{remark}

\section{The sufficiency trap}
\label{sec:trap}

Let $K_3$ denote the statement that every Kakeya set in $\mathbb{R}^3$ has
Hausdorff and Minkowski dimension three---now a theorem
\cite{WangZahl2025Kakeya}.

\begin{theorem}[Sufficiency trap]
\label{thm:trap}
\statusbox{\PeAIceFormal.}
Suppose one asserts the implication
\begin{equation}
  \label{eq:S}
  K_3 \;\Longrightarrow\; \RH.
\end{equation}
Since $K_3$ is a theorem, modus ponens would yield a proof of $\RH$.
Therefore either
\begin{enumerate}[label=(\roman*)]
  \item the implication \eqref{eq:S} is unproven and at least as hard as $\RH$, or
  \item asserting \eqref{eq:S} is asserting $\RH$ under a geometric label.
\end{enumerate}
In neither case does reading (S) carry information distinct from $\RH$ itself.
\end{theorem}

\begin{proof}
Immediate from modus ponens and the theoremhood of $K_3$.
If the implication remains unproved, its proof burden is at least that of $\RH$;
if it is taken as established, $\RH$ is established.
\end{proof}

\begin{corollary}
\label{cor:S-closed}
Reading (S) is \PeAIceClosedNeg\ as a research packaging of $\RH$.
Any non-trivial content of the Coleman Conjecture requires reading (N).
\end{corollary}

\begin{remark}
The trap is robust under alternate antecedents.
If the antecedent is still open (e.g.\ a Kakeya maximal-function conjecture, or
the set conjecture in dimensions $n\ge 4$), then ``antecedent $\Rightarrow\RH$''
is an unproven implication between hard problems, not a reduction that clarifies
either side.
\end{remark}

\section{The dependency tower and the Lindel\"of ceiling}
\label{sec:tower}

There is a genuine, citable tower in which Kakeya-type geometric control is
methodologically prior to analytic control of $\zeta$:
\begin{equation}
\label{eq:tower}
\begin{aligned}
&\text{incidence / tube geometry}
  \longrightarrow
\text{Kakeya / Besicovitch maximal phenomena}
\\
&\qquad\longrightarrow
\text{Fourier restriction / }\ell^2\text{-decoupling}
\\
&\qquad\longrightarrow
\text{Dirichlet-polynomial large values / exponential sums}
\\
&\qquad\longrightarrow
\text{zeta growth (subconvexity) and zero-density }\NsigT
\\
&\qquad\longrightarrow
\text{zero-free regions (partial)}.
\end{aligned}
\end{equation}
The rungs are established mathematics (\PeAIceKnown / theorem-background):
restriction interacts with Kakeya-type geometry \cite{Bourgain1991};
decoupling, built on wave-packet/tube organization, settled the main conjecture
in Vinogradov's mean value theorem \cite{BourgainDemeterGuth2016};
Guth--Maynard obtained new large-value estimates for Dirichlet polynomials and
deduced
\begin{equation}
  \label{eq:GM}
  \NsigT \;\ll\; T^{30(1-\sigma)/13+o(1)},
\end{equation}
a substantial improvement on classical Ingham-type density technology
\cite{GuthMaynard2024,Ingham1940}.

\begin{proposition}[Bound-type ceiling]
\label{prop:ceiling}
\statusbox{\PeAIceKnown\ for the forward implications; reverse open.}
Every rung of \eqref{eq:tower} produces \emph{bounds}, not exact zero locations.
In particular one has the standard forward chain
\begin{equation}
  \RH \;\Longrightarrow\; \LH \;\Longrightarrow\; \text{Density Hypothesis},
\end{equation}
where $\LH$ denotes the Lindel\"of hypothesis.
The reverse implications, including $\LH\Rightarrow\RH$, remain \PeAIceOpen.
\end{proposition}

\begin{remark}
Thus Kakeya-type geometry is vindicated as a \emph{necessary methodological
antecedent} of modern bound-type control of $\zeta$, and refuted as a
\emph{sufficient condition} for the critical line.
The meta-claim that \emph{no} bound-type method can reach $\RH$ is expert
expectation, not a theorem; we record it as \PeAIceAnalogy/\PeAIceOpen\
(Open Problem~\ref{op:ceiling}).
\end{remark}

\section{Invariant form and faithfulness}
\label{sec:invariant}

\begin{definition}[Shared invariant $\kappa$, schematic]
\label{def:kappa}
\statusbox{\PeAIceProposed.}
Let $\kappa$ denote a real parameter intended to encode simultaneously
\begin{enumerate}[label=(\roman*)]
  \item an optimal exponent in a Kakeya / restriction / maximal inequality, and
  \item a zero-density exponent $A$ in a bound of the form
  $\NsigT\ll T^{A(1-\sigma)+\varepsilon}$.
\end{enumerate}
A precise common normalization is not fixed here (Open Problem~\ref{op:kappa}).
\end{definition}

\begin{conjecture}[Coleman Conjecture, invariant form]
\label{conj:invariant}
\statusbox{\PeAIceProposed.}
There exists a normalization of $\kappa$ and an extremal value $\kappa^*$ such that
\begin{equation}
  \kappa=\kappa^* \quad\Longleftrightarrow\quad \RH.
\end{equation}
\end{conjecture}

\begin{proposition}[Forward direction, largely known]
\label{prop:forward}
\statusbox{\PeAIceKnown\ under standard consequences of $\RH$.}
Assuming $\RH$, one has $\NsigT=0$ for $\sigma>\tfrac12$ and Lindel\"of-type
bounds on the critical line; the tower exponents in Definition~\ref{def:kappa}
collapse to their extremal values, so $\kappa=\kappa^*$ holds for the naive
bound-type normalizations.
\end{proposition}

\begin{conjecture}[Reverse direction --- open core]
\label{conj:reverse}
\statusbox{\PeAIceOpen\ (at least as hard as $\RH$).}
$\kappa=\kappa^*\Rightarrow\RH$.
\end{conjecture}

\begin{proposition}[Faithfulness obstruction]
\label{prop:faithful}
\statusbox{\PeAIceFormal\ (conditional deduction).}
If $\kappa$ is identified with a pure harmonic-analytic / zero-density exponent
from Section~\ref{sec:tower}, then $\kappa=\kappa^*$ is Lindel\"of-scale
saturation (or a density-hypothesis extremal), and Conjecture~\ref{conj:reverse}
becomes an instance of ``bound-type extremality $\Rightarrow\RH$'', which is not
granted by known mathematics and is widely expected to fail for density proxies
alone.
Therefore a \emph{faithful} $\kappa$---one for which Conjecture~\ref{conj:invariant}
can hold non-vacuously---must encode strictly more than any bound-type
exponent: it must see exact zero \emph{locations}, not merely density.
\end{proposition}

\begin{proof}
Under the identification with tower exponents, extremality is a bound-type
statement.
The reverse implication of Conjecture~\ref{conj:reverse} then asserts that a
bound-type saturation yields $\RH$.
No such implication is known; the standard forward chain of
Proposition~\ref{prop:ceiling} does not reverse.
Hence faithfulness requires surplus structure beyond density exponents.
\end{proof}

\begin{remark}
Read invariantly, the Coleman Conjecture is therefore the bet that a faithful
Kakeya$\leftrightarrow\RH$ invariant \emph{exists}
(Open Problem~\ref{op:kappa}).
A conjecture may bet against prevailing expectation; the research-state grammar
requires the bet to be stated as such, not as a theorem.
\end{remark}

\section{Construction form and closed corridors}
\label{sec:construction}

A determinant identity of the shape
\begin{equation}
  \label{eq:det}
  \det_\zeta\bigl(B-z\bigr) \;=\; C\cdot \Xi(z)
\end{equation}
pins the zeros of the completed xi-function $\Xi$ exactly, not merely their
density---precisely the surplus demanded by Proposition~\ref{prop:faithful}.

\begin{conjecture}[Coleman Conjecture, construction form]
\label{conj:construction}
\statusbox{\PeAIceProposed.}
Any self-adjoint (or equivalently positivity-compatible) operator $B$ for which
a $\zeta$-regularized determinant identity \eqref{eq:det} holds must be built on
Kakeya-type incidence data (tube/shading geometry, multi-scale directional
organization).
In that sense Kakeya construction is a necessary antecedent of an
$\RH$-realizing spectral object.
\end{conjecture}

\subsection{Square-difference corridor}

A natural first attempt couples a thermal quadratic operator on a
$\Phi$-weighted $\ell^2$ space by a square-difference kernel
\begin{equation}
  \label{eq:Ksig}
  K_\sigma(m,n) \;=\; |m^2-n^2|^{-\sigma}\qquad(m\neq n),\qquad K_\sigma(n,n)=0.
\end{equation}
Program-level analysis of the Hilbert--Schmidt / determinant corridor for
\eqref{eq:Ksig} yields a closed-negative verdict for determinant-level
identification with $\Xi$: the kernel is Hilbert--Schmidt for $\sigma>\tfrac12$,
but singular-value / order-genus / counting constraints obstruct a
Riemann--von~Mangoldt density, and relatively compact perturbations of a
$D_1^2$-type operator preserve a counting incompatible with
$N(T)\sim (T/2\pi)\log T$
(program ledger: PeAIce DDATL / Claude~V6 wall registry; cf.\ also the
L2-5 counting mismatch discussion in the companion operator notes).

\begin{proposition}[Square-difference realization]
\label{prop:closed-neg}
\statusbox{\PeAIceClosedNeg\ (program corridor; not a claim about all operators).}
Within the PeAIce operator program, the pure square-difference coupling
\eqref{eq:Ksig} is closed as a route to \eqref{eq:det} with Riemann--von~Mangoldt
spectral counting.
\end{proposition}

\subsection{Prime-carrying requirement}

\begin{remark}[Surviving construction constraint]
\statusbox{\PeAIceProposed/\PeAIceOpen.}
Any surviving realization of \eqref{eq:det} is expected to be
\emph{prime-carrying} in the sense that prime lengths and von~Mangoldt-type
weights appear in the kernel or Hamiltonian data, with archimedean structure
producing
\begin{equation}
  N(T)\;\sim\; \frac{T}{2\pi}\log\frac{T}{2\pi},
\end{equation}
and with self-adjointness (or an equivalent positivity certificate).
Candidate external skeletons include adelic/trace-formula programs and
canonical-system / Krein constructions; converting those into a completed
identity \eqref{eq:det} remains \PeAIceOpen.
\end{remark}

\begin{remark}[Companion geometric notes]
Two companion notes in the same program formalize geometric layers used above
the slogan level:
a residual-overlap program after structured Kakeya factorization in
$\mathbb{R}^3$ (ARX-001 / Grain Zero), and a typed incidence object with a
non-inference firewall between overlap statistics and Hilbert-space placement
(ARX-002 / KNS(LB)).
Neither note proves $\RH$ or Conjecture~\ref{conj:cc-n}.
\end{remark}

\section{Open problems and falsifiers}
\label{sec:open}

\begin{enumerate}[label=\textbf{OP\arabic*.},leftmargin=*]
  \item \label{op:kappa}
  \textbf{Existence of a faithful $\kappa$.}
  \PeAIceOpen.
  Specify a single invariant bridging Kakeya/restriction data and exact zero
  location; by Proposition~\ref{prop:faithful} it must exceed every bound-type
  exponent.
  \item
  \textbf{Reverse implication.}
  \PeAIceOpen\ (RH-hard).
  Prove $\kappa=\kappa^*\Rightarrow\RH$ for a faithful $\kappa$.
  \item
  \textbf{Prime-carrying realization.}
  \PeAIceOpen\ (RH-hard).
  Construct a self-adjoint prime-carrying operator and prove \eqref{eq:det}
  with Riemann--von~Mangoldt counting.
  \item \label{op:ceiling}
  \textbf{Ceiling as theorem.}
  \PeAIceOpen.
  Make rigorous the extent to which bound-type statements in the lineage
  \eqref{eq:tower} can or cannot imply $\RH$.
\end{enumerate}

\paragraph{Falsifiers.}
\begin{enumerate}[label=\textbf{F\arabic*.},leftmargin=*]
  \item A $\kappa$-extremal configuration coexisting with a zero off the critical
  line falsifies the invariant form as stated.
  \item A self-adjoint realization of $\Xi$ with no Kakeya-type incidence data
  falsifies the necessity claim in the construction form.
  \item A proof that bound-type methods past Lindel\"of reach $\RH$ forces
  revision of the Kakeya-root framing of reading (N).
  \item A proof of $\LH\Rightarrow\RH$ would make a naive density $\kappa$
  faithful, collapsing Proposition~\ref{prop:faithful}---an independent
  revolution not specific to $\CC$.
\end{enumerate}

\section{Status summary}

\begin{center}
\begin{tabular}{@{}p{0.48\linewidth}p{0.22\linewidth}p{0.22\linewidth}@{}}
\hline
Statement & Reading & Status \\
\hline
Kakeya set theorem in $\mathbb{R}^3$ & background & \PeAIceKnown \\
Dependency tower rungs & background & \PeAIceKnown \\
$\CC$ as $K_3\Rightarrow\RH$ & (S) & \PeAIceClosedNeg \\
$\RH\Rightarrow\LH\Rightarrow$ Density Hyp. & ceiling & \PeAIceKnown \\
$\LH\Rightarrow\RH$ & ceiling & \PeAIceOpen \\
Guth--Maynard density bound & ceiling & \PeAIceKnown \\
Invariant form, forward & (N) & \PeAIceKnown \\
Invariant form, reverse & (N) & \PeAIceOpen \\
Faithfulness obstruction & deduction & \PeAIceFormal \\
Existence of faithful $\kappa$ & (N) & \PeAIceOpen \\
Construction form & (N) & \PeAIceProposed \\
Square-difference det corridor & construction & \PeAIceClosedNeg \\
Prime-carrying route & construction & \PeAIceOpen \\
$\RH$ and $\CC$ & --- & \PeAIceOpen \\
\hline
\end{tabular}
\end{center}

\section*{Acknowledgments}
This note formalizes an open research conjecture.
It does not claim a proof of the Riemann hypothesis.
Research engineering assistance was used in preparing program materials;
mathematical claims and any arXiv submission decision remain the author's
responsibility.
No model, webpage, or numerical run is sovereign over the mathematics.

\bibliographystyle{amsplain}
\bibliography{../../bibliography/peaice-arxiv}

\end{document}
